\section{Introduction}
\label{sec:introduction}

Neural networks are increasingly used as differentiable surrogates for
simulators that are too expensive to place directly inside design, estimation,
and control loops.  Once trained, a surrogate is often treated as an ordinary
smooth mapping and embedded in a nonlinear program (NLP).  This pattern appears
in process design, security-constrained power flow, adversarial verification,
and neural-operator inversion
\citep{yang2021modeling,kilwein2023feasibility,bugosen2023flowsheet,li2021fno}.
Modern automatic differentiation (AD) evaluates network
oracles efficiently, while mature primal--dual methods enforce nonlinear
constraints and report stationarity, feasibility, and complementarity
residuals.

The computational consequence is less clear.  A neural-surrogate NLP joins a
deep-learning runtime to a sparse optimization runtime.  Its wall time depends
on at least four coupled choices: (i) which derivative sparsity is disclosed to
the solver; (ii) whether curvature is exact, Gauss--Newton, or quasi-Newton;
(iii) whether framework boundaries stage tensors through host memory; and
(iv) whether the KKT system is assembled and factorized on the CPU or GPU.
Calling the network on a GPU resolves none of the other three choices.

Recent work has established full-space, reduced-space, and gray-box
formulations for neural predictors \citep{ceccon2022omlt,dowson2025mathoptai,
parker2024formulations}.  Most closely, \citet{parker2025gpu} move PyTorch
network derivatives to an A100 while Ipopt and its linear algebra remain on the
CPU.  In parallel, MadNLP and ExaModels demonstrate GPU-resident condensed KKT
methods for large algebraic NLPs
\citep{shin2024opf,pacaud2024dynamic,pacaud2024condensed}.  Together these
results motivate connecting a neural AD runtime to a GPU NLP solver.  Yet an end-to-end
implementation alone would answer only whether this is possible, not when it
is useful.  Small KKT systems cannot amortize GPU setup; a dense derivative
declaration can manufacture an artificial GPU crossover; and avoiding Hessian
callbacks can increase the number of expensive factorizations by an order of
magnitude.

Speed is not the only failure mode.  A KKT point is certified for
the \emph{surrogate} NLP, not for the physical system.  Optimizing a learned
model changes its input distribution and can amplify small prediction errors,
a central concern in offline model-based optimization
\citep{trabucco2021coms,trabucco2022designbench}.  For an inverse problem, a
method that achieves a lower surrogate residual can therefore produce a worse
simulator residual.  Systems benchmarks that omit independent replay can rank
methods precisely while ranking the wrong quantity.

This paper asks four questions:

\begin{enumerate}
  \item When does GPU-resident KKT linear algebra improve time to a
  solver-independent KKT target beyond GPU-accelerated neural AD alone?
  \item Can transfer removal matter after derivative sparsity and cold
  specialization are accounted for?
  \item How do exact, Gauss--Newton, and limited-memory curvature trade time,
  robustness, memory, and local solution quality, and does network topology
  matter after parameter count is fixed?
  \item Do rankings under the optimized surrogate agree with rankings after
  independent simulator replay?
\end{enumerate}

We answer these questions with \method, a JAX--MadNLP interface and an
experiment protocol constructed to vary neural parameter count, decision
dimension, and constraint/output dimension separately.  Our evidence combines
a planted smooth family, the public adversarial-MNIST problem used by the
closest GPU study, all three released neural-surrogate PFR-NMPC cases of
\citet{elorza2025nmpc}, and elliptic Darcy inversion with both a public FNO
split and a fully replayable finite-difference split.

Our contributions are:

\begin{itemize}
  \item \textbf{An audited end-to-end interface.}  DLPack shares CUDA buffers
  between CUDA.jl and JAX, derivative values are copied device-to-device into
  solver-owned storage, and evaluators declare sparse Jacobian and
  Lagrangian-Hessian coordinates.  Pointer identity, mutation, array layout,
  forward parity, and directional derivatives are separately tested.

  \item \textbf{A deployment regime map.}  We separate neural AD placement,
  host staging, KKT placement, curvature, sequence initialization, and cold
  specialization.  Across three published NMPC scales, correct sparsity first
  favors CPU KKT and a larger case then favors GPU KKT by
  \PFRCaseThreeSpeedup$\times$; both conclusions
  disappear in a one-dimensional ``GPU versus CPU'' label.

  \item \textbf{Failure-aware curvature evidence.}  A policy fixed on pilot
  cases runs L-BFGS first and invokes exact curvature only after a common-KKT
  failure.  On disjoint Darcy instances it obtains
  \DarcyRouteTotal/\DarcyRouteTotal{} successes with
  \DarcyRouteOverhead\% overhead over L-BFGS alone.  Other tasks show why this
  is a route rather than a universal L-BFGS recommendation.

  \item \textbf{Optimization-aware validity checks.}  Held-out targets,
  shared multi-starts, representability floors, raw residual audits,
  mechanistic closed-loop replay, and sparse PDE replay are part of the
  benchmark definition.  Failed runs and correctness-based exclusions remain
  documented in the supplementary material.  The
  replay study reveals that surrogate and simulator method rankings agree only
  \ReplayAgreement/\ReplayGroups{} times.
\end{itemize}

We make no new-formulation claim.  The empirical ordering is \emph{structure
before curvature, curvature before hardware, and simulator validation before
application claims}.  It turns an implementation
choice into a falsifiable experimental protocol and identifies which
accelerator regime is actually being measured.
